\documentclass{article}
\usepackage[utf8]{inputenc}

%%FOR CODES
\usepackage{listings,lstautogobble}
\usepackage{xcolor}
\usepackage{caption}
\usepackage{graphicx}

%%COLORES
\definecolor{realylightgray}{gray}{0.95}
\definecolor{RoyalBlue}{cmyk}{1, 0.50, 0, 0}
\definecolor{capri}{rgb}{0.0, 0.75, 1.0}
\definecolor{deepskyblue}{rgb}{0.0, 0.75, 1.0}
\definecolor{cobalt}{rgb}{0.0, 0.28, 0.67}
\definecolor{codegreen}{rgb}{0,0.6,0}
\definecolor{amber}{rgb}{1.0, 0.49, 0.0}
\definecolor{darkorchid}{rgb}{0.6, 0.2, 0.8}

\lstset{
	showstringspaces=false,
	showtabs=false,
	showspaces=false,
	stringstyle=\ttfamily,
	frameround=ffff,
	frame=single,
	framextopmargin=4pt,
	framexbottommargin=4pt,
	rulecolor=\color{darkgray},
	gobble=12, %align text inside box left <-> right
	belowskip=2em,
	keepspaces=true,
	escapeinside={<@}{@>},
	columns=fullflexible
}

%%template for commands
\lstdefinestyle{commands}{
	basicstyle=\scriptsize\ttfamily,
	commentstyle=\ttfamily\itshape\color{gray},
	backgroundcolor=\color{realylightgray},
	keywords=[1]{elbrus@server, root@server},
	keywordstyle=[1]\color{cobalt},
	numbers=none,
	breaklines=false
}

%%template for files
\lstdefinestyle{files}{
	basicstyle=\ttfamily\footnotesize,
	commentstyle=\color{codegreen},
	backgroundcolor=\color{realylightgray},    
	deletekeywords=[1]{elbrus@server, root@server},                                       
	numbers=left,       
	numbersep=5pt,
	breaklines=true,                                                 
}

\captionsetup[lstlisting]{singlelinecheck=false, } %to align caption left
\lstset{style=commands}

\begin{document}
	\begin{titlepage}
		\begin{center}
			\vspace*{1cm}
			
			\textbf{ELBRUS ANALYTICS - Bereitstellungshandbuch}
			
			\vspace{0.5cm}
			All about your Network
			
			\vspace{0.4cm}
			
			\includegraphics[width=0.2\textwidth]{logo.png}
			
			\vspace{0.6cm}
			\textbf{ELBRUS TEAM}
			
			\vfill
			
			Ein Handbuch zum Aufsetzten der \\
			Diplomarbeit Elbrus Analytics
			
			\vspace{0.8cm}
			
			\includegraphics[width=0.4\textwidth]{htl.png}
			
			Netzwerktechnik\\
			HTL Rennweg\\
			Österreich \\
			August 21, 2022
			
		\end{center}
	\end{titlepage}
	%\title{Elbrus Analytics - Bereitstellungshandbuch}
	%\author{ELBRUS TEAM}
	%\date{\today}
	%\maketitle
	\newpage
	Installationsanleitung für \textbf{Almalinux Skytiger (8.6)} bei abweichenden Versionen kann es zu Problemen bei der Installation kommen. \\ \\
	Bekanntes Problem: Beim herauskopieren von Befehlen wird das Apostrophe Zeichen falsch kopiert und führt zu Eingabe Störungen. Lösung: Apostrophe Zeichen des kopierten Befehls händisch im Terminal mit Apostrophe Zeichen austauschen. 
	
	\section[Server Infrastruktur]{Server Infrastruktur}
	\subsection{SSH-Zugriff vorbereiten}
	
	\lstset{style=commands}
	\begin{lstlisting}[caption={Updaten vorhandener Packages.}]
		root@server:~$ yum update -y
	\end{lstlisting}
	
	\lstset{style=commands}
	\begin{lstlisting}[caption={Installieren des 'ssh' Packages.}]
		root@server:~$ yum install -y openssh-server
	\end{lstlisting}

	\lstset{style=commands}
	\begin{lstlisting}[caption={Starten des 'sshd' Services.}]
		root@server:~$ systemctl start sshd
	\end{lstlisting}

	\lstset{style=commands}
	\begin{lstlisting}[caption={Aktivieren des 'sshd' Services.}]
		root@server:~$ systemctl enable sshd
	\end{lstlisting}

	\lstset{style=commands}
	\begin{lstlisting}[caption={Anlegen des Users Elbrus.}]
		root@server:~$ useradd elbrus
	\end{lstlisting}

	\lstset{style=commands}
	\begin{lstlisting}[caption={Hinzufügen des Users Elbrus zu der Gruppe 'wheel'.}]
		root@server:~$ usermod -aG wheel elbrus
	\end{lstlisting}
	\newpage
	
	\lstset{style=commands}
	\begin{lstlisting}[caption={Ändern des Passwords für den User Elbrus.}]
		root@server:~$ passwd elbrus
		Changing password for user elbrus.
		New password:
		Retype new password:
		passwd: all authentication tokens updated successfully
		root@server:~$
	\end{lstlisting}

	\lstset{style=commands}
	\begin{lstlisting}[caption={Wechseln zu User elbrus.}]
		root@server:~$ su elbrus
	\end{lstlisting}
	
	\subsection{Initiale Server Konfiguration}

	\lstset{style=commands}
	\begin{lstlisting}[caption={Setzen der Zeitzone auf 'Europa/Wien'.}]
		elbrus@server:~$ sudo timedatectl set-timezone Europe/Vienna
	\end{lstlisting}

	\lstset{style=commands}
	\begin{lstlisting}[caption={Installieren von dem 'firewalld' Service.}]
		elbrus@server:~$ sudo dnf install firewalld
	\end{lstlisting}
	\newpage
	
	\section{Voll automatisierte Installation}
	bei Verwendung der voll automatisierten Installation kann das restliche Bereitstellungshandbuch übersprungen werden. Falls eine manuelle Installation bevorzugt wird kann dieser Abschnitt übersprungen werden.
	
	\lstset{style=commands}
	\begin{lstlisting}[caption={Installieren des 'generate-keys' Script.}]
		elbrus@server:~$ sudo curl -o generate-keys.sh \
		https://directory.anzola.at/Elbrus/generate-keys.sh
	\end{lstlisting}

	\lstset{style=commands}
	\begin{lstlisting}[caption={Ausführen des 'generate-keys' Script.}]
		elbrus@server:~$ sudo bash generate-keys.sh
		Welchome to the setup of the 'elbrus analytics software'? (y/n) y
		
		starting to create ssh-keys to clone the github repositorys.
		
		Where should the ssh keys be stored (dir) [/var/elbrus/shared/key]:
		Please confirm the key folder '/var/elbrus/shared/key'. (y/n/exit) y
		
		Creating ssh-keys:
		[##########]
		Please send the generated public keys to keys@Elbrus-Analytics.at.
		elbrus@server:~$
	\end{lstlisting}
	
	Bevor mit der Installation fortgefahren werden kann, müssen die soeben angelegten SSH-Keys an "keys@Elbrus-Analytics.at" gesendet werden. Bitte beachten Sie, dass Ihre Email-Adresse Sie als berechtigten Nutzer ausweist.
	\newpage
	
	\lstset{style=commands}
	\begin{lstlisting}[caption={Installieren des 'install' Script.}]
		elbrus@server:~$ sudo curl -o install.sh https://directory.anzola.at/Elbrus/install.sh
	\end{lstlisting}

	\begin{lstlisting}[caption={Ausführen des 'install' Script.}]
		elbrus@server:~$ sudo bash install.sh
		Do you want to setup the 'elbrus analytics software'? (y/n) y
		
		Where are the previously generated ssh keys stored (dir) [/var/elbrus/shared/key]:
		Please confirm the key folder '/var/elbrus/shared/key'. (y/n/exit) y
		
		info: the 'elbrus analytics software' will be installed in '/var/elbrus'.
		info: downloading git.
		info: downloaded git.
		
		---------------------------------------------
		downloading the 'elbrus analytics software'
		---------------------------------------------
		
		info: checking downloaded directorys:
		1
		<@\textcolor{codegreen}{Successfully}@> installed the 'database' software.
		<@\textcolor{codegreen}{Successfully}@> installed the 'tabby' software.
		<@\textcolor{codegreen}{Successfully}@> installed the 'snmp-manager' software.
		<@\textcolor{codegreen}{Successfully}@> installed the 'ssh-manager' software.
		<@\textcolor{codegreen}{Successfully}@> installed the 'uptime-monitor' software.
		<@\textcolor{codegreen}{Successfully}@> installed the 'geo-session-finder' software.
		<@\textcolor{codegreen}{Successfully}@> installed the 'office365-analyzer' software.
		<@\textcolor{codegreen}{Successfully}@> installed the 'api' software.
		<@\textcolor{codegreen}{Successfully}@> installed the 'webinterface' software.
		
		...
		elbrus@server:~$
	\end{lstlisting}
	Im Falle von Unklarheiten bei Verwendung des 'install' Scripts kann bei jeder Teilinstallation im Bereitstellungshandbuch unter dem Abschnitt '1 - Mit Setup Script' der jeweilige Teilinstallation nachgeschlagen werden. Falls dies auch nicht behilflich ist überspringen Sie die jeweilige Teilinstallation und installieren Sie die jeweilige Teilinstallation manuell nach.
	\newline \newline \newline \newline \newline \newline \newline \newline \newline \newline \newline
	\textcolor{red}{Achtung!} Vergessen sie nicht in ihren Firewall Einstellungen die Ports für die API sowie das Webinteface frei zu geben.
	\newpage
	
	\subsection{Software Abhängigkeiten}
	
	\subsubsection{Node.Js}
	
	\lstset{style=commands}
	\begin{lstlisting}[caption={Installieren des Framworks 'Node.Js'.}]
		elbrus@server:~$ sudo dnf -y module install nodejs:12
	\end{lstlisting}

	\subsubsection{Nginx}
	
	\lstset{style=commands}
	\begin{lstlisting}[caption={Installieren der Webserver-Software 'Nginx'.}]
		elbrus@server:~$ sudo dnf -y install nginx
	\end{lstlisting}

	\subsubsection{Pm2}
		\lstset{style=commands}
	\begin{lstlisting}[caption={Installieren von 'pm2'.}]
		elbrus@server:/var/elbrus$ sudo npm install -g pm2
		elbrus@server:/var/elbrus$ sudo npm install --save luxon
	\end{lstlisting}
	
	\subsubsection{Ablagestruktur}
	
	\lstset{style=commands}
	\begin{lstlisting}[caption={Anlegen der Verzeichnissstruktur.}]
		elbrus@server:~$ sudo mkdir -p /var/elbrus/shared/keys
		elbrus@server:~$ sudo chown -R elbrus:elbrus /var/elbrus
		elbrus@server:~$ cd /var/elbrus
		elbrus@server:/var/elbrus$ chmod -R 777 /var/elbrus/shared
	\end{lstlisting}
	\newpage

	% -------- Git Repos clonen
	\subsection{Git}
	
	\lstset{style=commands}
	\begin{lstlisting}[caption={Installieren von dem VCS 'git'.}]
		elbrus@server:~$ sudo yum install -y git
	\end{lstlisting}
	
	\lstset{style=commands}
	\subsubsection{Git - Erstellen der SSH-Keys}
	\begin{lstlisting}[caption={Wechseln des Verzeichnisses.}]
		elbrus@server:~$ cd /var/elbrus/shared/keys
	\end{lstlisting}
	
	\lstset{style=commands}
	\begin{lstlisting}[caption={Erstellen des SSH-keys der für das Herunterladen der 'Database' benötigt wird.}]
		elbrus@server:/var/elbrus/keys$ ssh-keygen -t rsa -b 2048 -f database_key -q -N ""
	\end{lstlisting}

	\lstset{style=commands}
	\begin{lstlisting}[caption={Erstellen des SSH-Keys der für das Herunterladen der Kernsoftware 'tabby' benötigt wird.}]
		elbrus@server:/var/elbrus/keys$ ssh-keygen -t rsa -b 2048 -f tabby_key -q -N ""
	\end{lstlisting}

	\lstset{style=commands}
	\begin{lstlisting}[caption={Erstellen des SSH-keys der für das Herunterladen des 'SNMP-Managers' benötigt wird.}]
		elbrus@server:/var/elbrus/keys$ ssh-keygen -t rsa -b 2048 -f snmp_manager_key -q -N ""
	\end{lstlisting}

	\lstset{style=commands}
	\begin{lstlisting}[caption={Erstellen des SSH-keys der für das Herunterladen des 'SSH-Managers' benötigt wird.}]
		elbrus@server:/var/elbrus/keys$ ssh-keygen -t rsa -b 2048 -f ssh_manager_key -q -N ""
	\end{lstlisting}
	
	\lstset{style=commands}
	\begin{lstlisting}[caption={Erstellen des SSH-keys der für das Herunterladen des 'Uptime-Monitors' benötigt wird.}]
		elbrus@server:/var/elbrus/keys$ ssh-keygen -t rsa -b 2048 -f uptime_monitor_key -q -N ""
	\end{lstlisting}
	
	\lstset{style=commands}
	\begin{lstlisting}[caption={Erstellen des SSH-keys der für das Herunterladen des 'geo session finders' benötigt wird.}]
		elbrus@server:/var/elbrus/keys$ ssh-keygen -t rsa -b 2048 \
		 -f geo_session_finder_key -q -N ""
	\end{lstlisting}

	\lstset{style=commands}
	\begin{lstlisting}[caption={Erstellen des SSH-keys der für das Herunterladen des 'office365 analyzers' benötigt wird.}]
		elbrus@server:/var/elbrus/keys$ ssh-keygen -t rsa -b 2048 \
		-f office365_analyzer_key -q -N ""
	\end{lstlisting}
	
	\lstset{style=commands}
	\begin{lstlisting}[caption={Erstellen des SSH-keys der für das Herunterladen der 'API' benötigt wird.}]
		elbrus@server:/var/elbrus/keys$ ssh-keygen -t rsa -b 2048 -f api_key -q -N ""
	\end{lstlisting}
	
	\lstset{style=commands}
	\begin{lstlisting}[caption={Erstellen des SSH-keys der für das Herunterladen des 'Webinterfaces' benötigt wird.}]
		elbrus@server:/var/elbrus/keys$ ssh-keygen -t rsa -b 2048 -f webinterface_key -q -N ""
	\end{lstlisting}

	Bevor mit der Installation vorgefahren werden kann müssen die soeben angelegten SSH-Keys an "keys@Elbrus-Analytics.at" gesendet werden. Bitte beachten Sie, dass Ihre Email-Adresse Sie als berechtigten Nutzer ausweist.
	
	\subsubsection{Git - Clonen der Software}

	\lstset{style=commands}
	\begin{lstlisting}[caption={Wechseln des Verzeichnisses.}]
		elbrus@server:~$ cd /var/elbrus
	\end{lstlisting}
	
	\lstset{style=commands}
	\begin{lstlisting}[caption={Clonen der Datenbank Software.}]
		elbrus@server:/var/elbrus$ git clone git@github.com:Elbrus-Analytics/database.git \
		--config core.sshCommand="ssh -i /var/elbrus/shared/keys/database_key"
	\end{lstlisting}

	\lstset{style=commands}
	\begin{lstlisting}[caption={Clonen der Kernsoftware 'tabby'.}]
		elbrus@server:/var/elbrus$ git clone git@github.com:Elbrus-Analytics/tabby.git \
		--config core.sshCommand="ssh -i /var/elbrus/shared/keys/tabby_key"
	\end{lstlisting}
	
	\lstset{style=commands}
	\begin{lstlisting}[caption={Clonen der 'SNMP-Manager' Software.}]
		elbrus@server:/var/elbrus$ git clone git@github.com:Elbrus-Analytics/snmp-manager.git \
		--config core.sshCommand="ssh -i /var/elbrus/shared/keys/snmp_manager_key"
	\end{lstlisting}

	\lstset{style=commands}
	\begin{lstlisting}[caption={Clonen der 'SSH-Manager' Software.}]
		elbrus@server:/var/elbrus$ git clone git@github.com:Elbrus-Analytics/ssh-manager.git \
		--config core.sshCommand="ssh -i /var/elbrus/shared/keys/ssh_manager_key"
	\end{lstlisting}
	
	\lstset{style=commands}
	\begin{lstlisting}[caption={Clonen der 'Uptime-Monitor' Software.}]
		elbrus@server:/var/elbrus$ git clone git@github.com:Elbrus-Analytics/uptime-monitor.git \
		--config core.sshCommand="ssh -i /var/elbrus/shared/keys/uptime_monitor_key"
	\end{lstlisting}

	\lstset{style=commands}
	\begin{lstlisting}[caption={Clonen der 'geo session finders' Software.}]
		elbrus@server:/var/elbrus/keys$  git clone git@github.com:Elbrus-Analytics/\
		geo-session-finder.git --config core.sshCommand="ssh -i \
		/var/elbrus/shared/keys/geo_session_finder_key"
	\end{lstlisting}

	\lstset{style=commands}
	\begin{lstlisting}[caption={Clonen der 'office365-analyzer' Software.}]
		elbrus@server:/var/elbrus/keys$ git clone git@github.com:Elbrus-Analytics/\
		office365-analyzer.git --config core.sshCommand="ssh -i \
		/var/elbrus/shared/keys/office365_analyzer_key"
	\end{lstlisting}
	
	\lstset{style=commands}
	\begin{lstlisting}[caption={Clonen der 'API' Software.}]
		elbrus@server:/var/elbrus$ git clone git@github.com:Elbrus-Analytics/api.git \
		--config core.sshCommand="ssh -i /var/elbrus/shared/keys/api_key"
	\end{lstlisting}

	\lstset{style=commands}
	\begin{lstlisting}[caption={Clonen der 'Webinterface' Software}]
		elbrus@server:/var/elbrus$ git clone git@github.com:Elbrus-Analytics/webinterface.git \
		--config core.sshCommand="ssh -i /var/elbrus/shared/keys/webinterface_key"
	\end{lstlisting}
	\newpage
	
	\subsection{Python}
	
	\lstset{style=commands}
	\begin{lstlisting}[caption={Installieren von benötigten Packeten und Abhängigkeiten.}]
		elbrus@server:~$ cd /var/elbrus
		elbrus@server:/var/elbrus$ sudo dnf install gcc openssl-devel bzip2-devel\
		libffi-devel zlib-devel wget make -y
	\end{lstlisting}
	
	\lstset{style=commands}
	\begin{lstlisting}[caption={Herunterladen der Source Datei.}]
		elbrus@server:/var/elbrus$ wget https://www.python.org/ftp/python/\
		3.10.0/Python-3.10.0.tar.xz
	\end{lstlisting}
	
	\lstset{style=commands}
	\begin{lstlisting}[caption={Extrahieren der installierten Datei.}]
		elbrus@server:/var/elbrus$ tar -xf Python-3.10.2.tar.xz
	\end{lstlisting}
	
	\lstset{style=commands}
	\begin{lstlisting}[caption={Wechseln zu source Verzeichniss. Und ausführen des Konfigurations Scripts.}]
		elbrus@server:/var/elbrus$ cd Python-3.10.0 && ./configure --enable-optimizations
	\end{lstlisting}

	\lstset{style=commands}
	\begin{lstlisting}[caption={Starten des build Prozesses.}]
		elbrus@server:/var/elbrus/Python-3.10.0$ cd make -j $(nproc)
	\end{lstlisting}
	
	\lstset{style=commands}
	\begin{lstlisting}[caption={Installieren von Python.}]
		elbrus@server:/var/elbrus/Python-3.10.0$ sudo make install
	\end{lstlisting}

	\lstset{style=commands}
	\begin{lstlisting}[caption={Löschen der kompremierten Python Datei.}]
		elbrus@server:/var/elbrus/Python-3.10.0$ cd .. && rm Python-3.10.0.tar.xz
	\end{lstlisting}

	\subsubsection{Upgrade von 'pip'}

	\lstset{style=commands}
	\begin{lstlisting}[caption={Upgraden von 'pip'.}]
		elbrus@server:~$ /usr/local/bin/python3.10 -m pip install --upgrade pip
	\end{lstlisting}

	\subsection{Rust}
	
	\lstset{style=commands}
	\begin{lstlisting}[caption={Installieren von GNU Compiler Collection.}]
		elbrus@server:~$ sudo dnf install gcc -y
	\end{lstlisting}
	
	\lstset{style=commands}
	\begin{lstlisting}[caption={Installieren von Rust.}]
		elbrus@server:~$ curl --proto '=https' --tlsv1.2 -sSf \
		https://sh.rustup.rs/ | sh
		 
		...
		 
		default host triple: x86_64-unknown-linux-gnu
		default toolchain: stable (default)
		profile: default
		modify PATH variable: yes
		
		1) Proceed with installation (default)
		2) Customize installation
		3) Cancel installation
		>1
		
		...
		
		<@\textcolor{codegreen}{stable-x86\_64-unknown-linux-gnu installed}@> - rustc 1.62.1 (e092d0b6b 2022-07-16)
		
		
		Rust is installed now. Great!
		
		To get started you may need to restart your current shell.
		This would reload your PATH environment variable to include
		Cargo's bin directory ($HOME/.cargo/bin).
		
		To configure your current shell, run:
		source "$HOME/.cargo/env"
		elbrus@server:~$
	\end{lstlisting}
	
	\lstset{style=commands}
	\begin{lstlisting}[caption={Laden der Variablen aus dem Terminal Profil.}]
		elbrus@server:~$ source ~/.profile
	\end{lstlisting}

	\lstset{style=commands}
	\begin{lstlisting}[caption={Hinzufügen des Befehls Cargo zu dem Pfad.}]
		elbrus@server:~$ source ~/.cargo/env
	\end{lstlisting}
	\newpage
	
	% --------- DATENBANK ---------
	\section{Datenbank}
	% --------- deployment guide for the database ---------
	\subsection[dependencies]{Voraussetzungen}
	
	\lstset{style=commands}
	\begin{lstlisting}[caption={Hinzufügen des PostgreSQL Drittanbieter-Repository, um die neuesten PostgreSQL-Pakete zu erhalten.}]
		elbrus@server:~$ sudo yum install\
		https://download.postgresql.org/pub/repos/yum/reporpms/\
		EL-$(rpm -E %{rhel})-x86_64/pgdg-redhat-repo-latest.noarch.rpm
	\end{lstlisting}
	
	\lstset{style=commands}
	\begin{lstlisting}[caption={Erstellen des Timescale repository.}]
		elbrus@server:~$ sudo tee /etc/yum.repos.d/timescale_timescaledb.repo <<EOL
		[timescale_timescaledb]
		name=timescale_timescaledb
		baseurl=https://packagecloud.io/timescale/timescaledb/el/$(rpm -E %{rhel})/$basearch
		repo_gpgcheck=1
		gpgcheck=0
		enabled=1
		gpgkey=https://packagecloud.io/timescale/timescaledb/gpgkey
		sslverify=1
		sslcacert=/etc/pki/tls/certs/ca-bundle.crt
		metadata_expire=300
		EOL
	\end{lstlisting}
	
	\lstset{style=commands}
	\begin{lstlisting}[caption={Updaten der lokalen Package-Liste.}]
		elbrus@server:~$ sudo yum update -y
	\end{lstlisting}
	
	\lstset{style=commands}
	\begin{lstlisting}[caption={Installieren von TimescaleDB.}]
		elbrus@server:~$ sudo dnf -qy module disable postgresql
		elbrus@server:~$ sudo dnf install postgresql14 postgresql14-server -y
		elbrus@server:~$ sudo dnf install timescaledb-2-postgresql-14 -y
	\end{lstlisting}
	\newpage
	
	\subsection[TimescaleDB konfigurieren]{Umgebung Konfigurieren}
	
	\lstset{style=commands}
	\begin{lstlisting}[caption={Initialisieren der Datenbank.}]
		elbrus@server:~$ sudo /usr/pgsql-14/bin/postgresql-14-setup initdb
	\end{lstlisting}
	
	\lstset{style=commands}
	\begin{lstlisting}[caption={Verknüpfen von 'postgresql' Serive Start mit Serverstart sowie den Service starten.}]
		elbrus@server:~$ sudo systemctl enable postgresql-14
		elbrus@server:~$ sudo systemctl start postgresql-14
	\end{lstlisting}
	
	\lstset{style=files}
	\begin{lstlisting}[caption={/var/lib/pgsql/14/data/postgresql.conf - Ändern der folgenden Zeilen.}, numbers=none]
		<@\textcolor{red}{- \#shared\_preload\_libraries = ''}@>
		<@\textcolor{codegreen}{+ shared\_preload\_libraries = 'timescaledb'}@>
		
		<@\textcolor{red}{- \#listen\_addresses = 'localhost'}@>
		<@\textcolor{codegreen}{+ listen\_addresses = '*'}@>
	\end{lstlisting}
	
	\lstset{style=files}
	\begin{lstlisting}[caption={/var/lib/pgsql/14/data/pg\_hba.conf - Ändern der folgenden Zeilen.}, numbers=none]
		# TYPE      DATABASE      USER      ADDRESS      METHOD
		<@\textcolor{codegreen}{+ host}@>      <@\textcolor{codegreen}{elbrus}@>        <@\textcolor{codegreen}{elbrus}@>    <@\textcolor{codegreen}{0.0.0.0/0}@>    <@\textcolor{codegreen}{trust}@>
	\end{lstlisting}
	
	\lstset{style=commands}
	\begin{lstlisting}[caption={Anpassen der Datenbank Einstellungen auf die Server Hardware.}]
		elbrus@server:~$ sudo timescaledb-tune --pg-config=/usr/\
		pgsql-14/bin/pg_config --yes
	\end{lstlisting}
	
	\lstset{style=commands}
	\begin{lstlisting}[caption={Neustarten des Services um Änderungen zu übernehmen.}]
		elbrus@server:~$ sudo systemctl restart postgresql-14
	\end{lstlisting}
	\newpage
	
	\subsection{Erstellen der Elbrus-Datenbank}
	
	\lstset{style=commands}
	\begin{lstlisting}[caption={Verbinden mit dem interaktiven Terminal von 'postgres'.}]
		elbrus@server:~$ sudo su postgres -c psql
	\end{lstlisting}
	Im folgenden Text sind markierte Abschnitte Variablen, welche im darunterliegen SQL gerändert werden können, was aus Sichertsgründen dringend empfohlen wird.
	\begin{enumerate}
		\item Die Datenbank elbrus anlegen
		\item Die Zeitzone auf Europe/Vienna setzen
		\item Den User elbrus mit dem Passwort elbrus123! anlegen
		\item Dem User alle rechte auf die voher erstellte Datenbank geben
	\end{enumerate}	
	
	\lstset{style=files}
	\begin{lstlisting}[caption={Auführen von SQL Befehlen.}, numbers=none]
		<@\textcolor{amber}{CREATE DATABASE}@> <@\textcolor{darkorchid}{ elbrus}@>;
		<@\textcolor{amber}{ALTER DATABASE}@> elbrus <@\textcolor{amber}{SET}@> timezone TO <@\textcolor{deepskyblue}{'Europe/Vienna'}@>;
		<@\textcolor{amber}{CREATE USER}@> <@\textcolor{darkorchid}{ elbrus}@> PASSWORD <@\textcolor{deepskyblue}{'elbrus123!'}@>;
		<@\textcolor{amber}{GRANT}@> ALL <@\textcolor{amber}{ON}@> DATABASE elbrus TO elbrus;
	\end{lstlisting}
	
	\begin{lstlisting}[caption={Wechseln zu erstellter Datenbank.}, numbers=none]
		\c elbrus
	\end{lstlisting}
	
	\begin{lstlisting}[caption={Hinzufügen der TimescaleDB Erweiterung.}, numbers=none]
		<@\textcolor{amber}{CREATE EXTENSION IF NOT EXISTS}@> timescaledb;
		exit
	\end{lstlisting}
	
	\lstset{style=commands}
	\begin{lstlisting}[caption={Anlegen der benötigten Tabellen duch das ausführen von 'init.sql'.}]
		elbrus@server:/var/elbrus$ psql -U elbrus -d elbrus -f \
		database/sql/init.sql
	\end{lstlisting}
	\newpage
	
	% ----------- Globale Konfiguration ---------
	\section{Globale Konfiguration}
	% --------- deployment guide for the global konfiguration ---------
	\lstset{style=files}
	\begin{lstlisting}[caption={Anhand von '.config.example' eigene '.config' Datei in \newline'/var/elbrus/shared' anlegen.}, language=bash]
		#database settings
		DB_HOST=localhost
		DB_PORT=5432
		DB_NAME=elbrus
		DB_USER=elbrus
		DB_PASSWORD=elbrus123!
		
		#paths
		PCAPFOLDER=/var/elbrus/shared/traces
		IMPORTERPATH=/var/elbrus/pcap-importer/pcap-importer
		REPORTERPATH=/var/elbrus/report-generator/src/main.py
	\end{lstlisting}
	
	\lstset{style=commands}
	\begin{lstlisting}
		elbrus@server:/var/elbrus$ sudo chown elbrus:elbrus /var/elbrus/shared/.config
		elbrus@server:/var/elbrus$ sudo chmod 776 /var/elbrus/shared/.config
	\end{lstlisting}
	\newpage
	
	% --------- TABBY ---------
	\section{Kernsoftware - Tabby}
	% --------- deployment guide for tabby ---------		
	\subsection{1 - Mit Setup Script}
	
	\lstset{style=commands}
	\begin{lstlisting}[caption={Ausführen des 'install.sh' Scripts.}]
		elbrus@server:~$ cd /var/elbrus
		elbrus@server:/var/elbrus$ sudo bash tabby/install.sh
		Do you want to proceed with setup of the 'tabby'? (y/n) y
		
		Where is the shared config stored [/var/elbrus/shared/.config]:
		Where should the tabby be stored [/var/elbrus/tabby]:
		Where should the log be stored (dir) [/var/elbrus/shared/log]:
		Where should the traces be stored (dir) [/var/elbrus/shared/traces]:
		
		Is the shared config stored at '/var/elbrus/shared/.config' ?
		Should the capture-device be stored at '/var/elbrus/tabby' ?
		Should the log be stored at '/var/elbrus/shared/log' ?
		Should the traces be stored at '/var/elbrus/shared/traces' ? (y/n/exit) y
		
		info: a listing of your interfaces follows
		
		...
		
		Which inteface should be captured (name)? eth0if81
		
		Should 'eth0if81' be captured? (y/n/exit) y
		
		info: the capturing interface can be changed in the environment file.
		
		info: building the software using rust
		
		<@\textcolor{codegreen}{Success!}@> built the software
		<@\textcolor{codegreen}{Success!}@> .env file was created
		<@\textcolor{codegreen}{Success!}@> systemd service was automatically deployed
		elbrus@server:/var/elbrus$
	\end{lstlisting}
	\newpage
	
	\subsection{2 - Ohne Setup Script}
	\subsubsection[file config]{Umgebung Konfigurieren}
	\lstset{style=files}
	\begin{lstlisting}[caption={Anhand von '.env.example' eigene '.env' Datei anlegen.}, language=bash]
		#global
		SHAREDCONFIG=/var/elbrus/shared/.config
		
		#paths
		PCAPFOLDER=/var/elbrus/shared/traces
		LOGFILEDIR=/var/elbrus/shared/log
		CAPTUREPATH=/var/elbrus/tabby/tabby
		
		#settings
		TIMEPERCAPTURE=600
		TIMEPERIMPORT=300
		TIMEPERREPORT=900
		MAXFILES=12
		INTERFACE=eth0if81
	\end{lstlisting}
	
	\lstset{style=commands}	
	\begin{lstlisting}[caption={Kompilieren der Kernsoftware.}]
		elbrus@server:/var/elbrus$ sudo /root/.cargo/bin/cargo build --release \
		--manifest-path="/var/elbrus/tabby/Cargo.toml"
	\end{lstlisting}
	
	\lstset{style=commands}
	\begin{lstlisting}[caption={Verschieben der kompilierten Software.}]
		elbrus@server:/var/elbrus$ mv tabby/target/release/tabby tabby/tabby
	\end{lstlisting}

	\lstset{style=commands}
	\begin{lstlisting}[caption={Verschieben der kompilierten Software.}]
		elbrus@server:/var/elbrus$ mv tabby/target/release/tabby tabby/tabby
	\end{lstlisting}

	\lstset{style=commands}
	\begin{lstlisting}[caption={Anlegen des Benutzers 'tabby'.}]
		elbrus@server:/var/elbrus$ sudo useradd -G elbrus tabby
	\end{lstlisting}

	\lstset{style=commands}
	\begin{lstlisting}[caption={Verändern der berechtigungen auf Ordner 'tabby'.}]
		elbrus@server:/var/elbrus$ sudo chown -R tabby:tabby tabby
	\end{lstlisting}
	\newpage

	\lstset{style=commands}
	\begin{lstlisting}[caption={Anlegen des Log-Verzeichnisses.}]
		elbrus@server:/var/elbrus$ mkdir -p /var/elbrus/shared/log
	\end{lstlisting}

	\lstset{style=commands}
	\begin{lstlisting}[caption={Verändern der Berechtigung das Log-Verzeichnisses.}]
		elbrus@server:/var/elbrus$ chmod 777 -R /var/elbrus/shared/log
	\end{lstlisting}

	\lstset{style=commands}
	\begin{lstlisting}[caption={Verändern der Berechtigung der kompilierten Software.}]
		elbrus@server:/var/elbrus$ sudo chmod 750 tabby/tabby
		elbrus@server:/var/elbrus$ sudo setcap cap_net_raw,cap_net_admin=eip tabby/tabby
	\end{lstlisting}

	\subsubsection[systemd service]{Der Systemd Service}
	
	\lstset{style=files}
	\begin{lstlisting}[caption={tabby.service.example - Die Variable 'WorkingDirectory' sowie die Variable 'User' anpassen.},language=bash ,keywords={WorkingDirectory, User}, keywordstyle=\color{red}, firstnumber=5]
		...
		#job is starting immediatly after the start action has been called
		Type=simple
		#the user to execute the script
		User=tabby
		#the working directory
		WorkingDirectory=/var/elbrus/tabby/
		#which script should be executed
		ExecStart=/bin/bash tabby.sh
		...
	\end{lstlisting}

	\lstset{style=files}
	\begin{lstlisting}[caption={tabby-error-handler.service.example - Die Variable 'WorkingDirectory' sowie die Variable 'User' anpassen.},language=bash ,keywords={WorkingDirectory, User}, keywordstyle=\color{red}, firstnumber=5]
		...
		[Service]
		Type=oneshot
		User=tabby
		WorkingDirectory=/var/elbrus/tabby/
		ExecStart=/bin/bash tabby-log.sh
		...
	\end{lstlisting}
	\newpage
	
	\lstset{style=commands}
	\begin{lstlisting}[caption={Kopieren des Serviceprogrammes.}]
		elbrus@server:/var/elbrus$ sudo cp tabby/tabby.service.example \
		/etc/systemd/system/tabby.service
	\end{lstlisting}
	
	\lstset{style=commands}
	\begin{lstlisting}[caption={Kopieren des Errorhandlers.}]
		elbrus@server:/var/elbrus$ sudo cp tabby/tabby-error-handler.service.example \
		/etc/systemd/system/tabby-error-handler.service
	\end{lstlisting}
	
	\lstset{style=commands}
	\begin{lstlisting}[caption={Neuladen des 'systemctl' Deamons.}]
		elbrus@server:/var/elbrus$ sudo systemctl daemon-reload
	\end{lstlisting}
	
	\lstset{style=commands}
	\begin{lstlisting}[caption={Aktivieren des Serviceprogrammes.}]
		elbrus@server:/var/elbrus$ sudo systemctl enable tabby.service
	\end{lstlisting}
	
	\lstset{style=commands}
	\begin{lstlisting}[caption={Starten des Serviceprogrammes.}]
		elbrus@server:/var/elbrus$ sudo systemctl start tabby.service
	\end{lstlisting}
	\newpage
	
	% --------- SNMP MANAGER ---------
	\section{SNMP Manager}
	% --------- deployment guide for the snmp manager ---------		
	\subsection{1 - Mit Setup Script}
	
	\lstset{style=commands}
	\begin{lstlisting}[caption={Ausführen des 'install.sh' Scripts.}]
		elbrus@server:~$ cd /var/elbrus
		elbrus@server:/var/elbrus$ sudo bash snmp-manager/src/install.sh 
		Do you want to proceed with setup of the 'snmp-manager'? (y/n) y
		
		Where should the log be stored (dir) [/var/elbrus/shared/log]: 
		Where should the 'snmp-manager' be stored [/var/elbrus/snmp-manager]:
		Where is the shared config stored [/var/elbrus/shared/.config]: 
		
		Should the log be stored at '/var/elbrus/shared/log' ?
		Should the 'snmp-manager' be stored at '/var/elbrus/snmp-manager'?
		Is the shared config stored at '/var/elbrus/shared/.config' ? (y/n/exit) y
		
		<@\textcolor{codegreen}{Success!}@> .env file was created
		<@\textcolor{codegreen}{Success!}@> systemd service was automatically deployed,
		<@\textcolor{codegreen}{Success!}@> installed requirements
		
		elbrus@server:/var/elbrus$
	\end{lstlisting}
	
	\subsection{2 - Ohne Setup Script}
	
	\subsubsection[file config]{Umgebung Konfigurieren}
	\lstset{style=files}
	\begin{lstlisting}[caption={Anhand von '.env.example' eigene '.env' Datei anlegen.}, language=bash, keywords={CONFIGPATH, MAINPATH}, keywordstyle=\color{red}]
		#global
		SHAREDCONFIG=/var/elbrus/shared/.config
		
		#paths
		LOGFILEDIR=/var/elbrus/shared/log
	\end{lstlisting}
	\newpage
	
	\subsubsection[systemd service]{Der Systemd Service}
	
	\lstset{style=files}
	\begin{lstlisting}[caption={snmp-manager.service.example - Die Variable 'WorkingDirectory' sowie die Variable 'User' anpassen.},language=bash ,keywords={WorkingDirectory, User}, keywordstyle=\color{red}, firstnumber=5]
		...
		#job is starting immediatly after the start action has been called
		Type=simple
		#the user to execute the script
		User=elbrus
		#the working directory
		WorkingDirectory=/var/elbrus/snmp-manager/src
		#which script should be executed
		ExecStart=/bin/bash elb-snmp-manager.sh
		...
	\end{lstlisting}
	
	\lstset{style=commands}
	\begin{lstlisting}[caption={Kopieren des Serviceprogrammes.}]
		elbrus@server:/var/elbrus$ sudo cp snmp-manager/src/snmp-manager.service\
		.example /etc/systemd/system/snmp-manager.service
	\end{lstlisting}
	
	\lstset{style=commands}
	\begin{lstlisting}[caption={Kopieren des Zeitplanungsprogrammes.}]
		elbrus@server:/var/elbrus$ sudo cp snmp-manager/src/snmp-manager-schedule.timer\
		.example /etc/systemd/system/snmp-manager-schedule.timer
	\end{lstlisting}
	
	\lstset{style=commands}
	\begin{lstlisting}[caption={Neuladen des 'systemctl' Deamons.}]
		elbrus@server:/var/elbrus$ sudo systemctl daemon-reload
	\end{lstlisting}
	
	\lstset{style=commands}
	\begin{lstlisting}[caption={Aktivieren des Serviceprogrammes.}]
		elbrus@server:/var/elbrus$ sudo systemctl enable snmp-manager.service
	\end{lstlisting}
	
	\lstset{style=commands}
	\begin{lstlisting}[caption={Aktivieren des Zeitplanungsprogrammes.}]
		elbrus@server:/var/elbrus$ sudo systemctl enable snmp-manager-schedule.timer
	\end{lstlisting}
	
	\lstset{style=commands}
	\begin{lstlisting}[caption={Starten des Zeitplanungsprogrammes.}]
		elbrus@server:/var/elbrus$ sudo systemctl start snmp-manager-schedule.timer
	\end{lstlisting}
	\newpage
	
	% --------- SSH MANAGER ---------
	\section{SSH Manager}
	% --------- deployment guide for the ssh manager ---------		
	\subsection[file config]{Umgebung Konfigurieren}
	\subsubsection{1 - Mit Setup Script}
	
	\lstset{style=commands}
	\begin{lstlisting}[caption={Ausführen des 'install.sh' Scripts.}]
		elbrus@server:~$ cd /var/elbrus
		elbrus@server:~$ bash ssh-manager/src/install.sh
		Do you want to proceed with the setup of the 'ssh-manager'? (y/n) y
		we will proceed
		
		Where do you want the ssh config replies to be stored (dir) [/var/elbrus/
		shared/ssh-configs]
		Where is the 'main.py' file stored [/var/elbrus/ssh-manager/src/main.py]
		Where is the shared config stored [/var/elbrus/shared/.config]
		
		Do you want to store the config files at '/var/elbrus/shared/ssh-configs'?
		Is your 'main.py' stored at '/var/elbrus/ssh-manager/src/main.py'?
		Is the shared config stored at '/var/elbrus/shared/.config'? (y/n/exit) y
		we will proceed
		
		#global
		SHAREDCONFIG=/var/elbrus/shared/.config
		
		#values regarding the jumpserver:
		#IP, PORT and USER values must be set!
		#depending on the usage you can set either:
		#   -PASS and KEYFILE: the keyfile is used, the pass is interpreted as the
		passphrase
		#   -only KEYFILE: the keyfile is used
		#   -only PASS: the password is used as is regular credentials
		JUMPSERVER_IP=
		JUMPSERVER_PORT=
		JUMPSERVER_USER=
		JUMPSERVER_PASS=
		SSH_KEYFILE=
		
		#paths
		CONFIGPATH=/var/elbrus/shared/ssh-configs
		MAINPATH=/var/elbrus/ssh-manager/src/main.py
		
		Do you want to run the setup script? (y/n/exit) y
		...
		Initialized empty Git repository in /var/elbrus/shared/ssh-configs/.git/
		info: created config folder!
		Install dependencies ...
		...
		Cleaning up...
		elbrus@server:/ssh-manager$
	\end{lstlisting}
	\newpage
	
	\lstset{style=files}
	\begin{lstlisting}[caption={Ergänzen der fehlenden Werten in '.env'.}, language=bash]
		#values regarding the jumpserver:
		#IP, PORT and USER values must be set!
		#depending on the usage you can set either:
		#   -PASS and KEYFILE: the keyfile is used, the pass is interpreted as the passphrase
		#   -only KEYFILE: the keyfile is used
		#   -only PASS: the password is used as is regular credentials
		JUMPSERVER_IP=
		JUMPSERVER_PORT=
		JUMPSERVER_USER=
		JUMPSERVER_PASS=
		SSH_KEYFILE=
	\end{lstlisting}
	
	\subsubsection{2 - Ohne Setup Script}
	
	\lstset{style=files}
	\begin{lstlisting}[caption={Anhand von '.env.example' eigene '.env' Datei anlegen.}, language=bash]
		#global
		SHAREDCONFIG=/var/elbrus/shared/.config
		
		#values regarding the jumpserver:
		#IP, PORT and USER values must be set!
		#depending on the usage you can set either:
		#   -PASS and KEYFILE: the keyfile is used, the pass is interpreted as the passphrase
		#   -only KEYFILE: the keyfile is used
		#   -only PASS: the password is used as is regular credentials
		JUMPSERVER_IP=
		JUMPSERVER_PORT=
		JUMPSERVER_USER=
		JUMPSERVER_PASS=
		SSH_KEYFILE=
		
		#paths
		CONFIGPATH=/var/elbrus/shared/ssh-configs
		MAINPATH=/var/elbrus/ssh-manager/src/main.py
	\end{lstlisting}
	
	\lstset{style=commands}	
	\begin{lstlisting}[caption={Ausführen des Scripts zur Initialisierung des VCS Verzeichnisses.}]
		elbrus@server:/var/elbrus/ssh-manager$ bash src/setup.sh
	\end{lstlisting}
	
	\lstset{style=commands}
	\begin{lstlisting}[caption={Installieren von fehlenden python3 Packages.}]
		elbrus@server:/var/elbrus/ssh-manager$ pip3 install -r requirements.txt
	\end{lstlisting}
	\newpage
	
	\subsection[systemd service]{Der Systemd Service}
	
	\lstset{style=files}
	\begin{lstlisting}[caption={ssh-manager.service.example - Die Variable 'WorkingDirectory' sowie die Variable 'User' anpassen.},language=bash ,keywords={WorkingDirectory, User}, keywordstyle=\color{red}, firstnumber=5]
		...
		#job is starting immediatly after the start action has been called
		Type=simple
		#the user to execute the script
		User=elbrus
		#the working directory
		WorkingDirectory=/var/elbrus/ssh-manager/src/
		#which script should be executed
		ExecStart=/bin/bash routine.sh
		...
	\end{lstlisting}
	
	\lstset{style=commands}
	\begin{lstlisting}[caption={Kopieren des Serviceprogrammes.}]
		elbrus@server:/var/elbrus$ sudo cp ssh-manager/src/ssh-manager.service.example \
		/etc/systemd/system/ssh-manager.service
	\end{lstlisting}
	
	\lstset{style=commands}
	\begin{lstlisting}[caption={Kopieren des Zeitplanungsprogrammes.}]
		elbrus@server:/var/elbrus$ sudo cp ssh-manager/src/ssh-manager-schedule.timer.example \
		/etc/systemd/system/ssh-manager-schedule.timer
	\end{lstlisting}
	
	\lstset{style=commands}
	\begin{lstlisting}[caption={Neuladen des 'systemctl' Deamons.}]
		elbrus@server:/var/elbrus$ sudo systemctl daemon-reload
	\end{lstlisting}
	
	\lstset{style=commands}
	\begin{lstlisting}[caption={Aktivieren des Serviceprogrammes.}]
		elbrus@server:/var/elbrus$ sudo systemctl enable ssh-manager.service
	\end{lstlisting}
	
	\lstset{style=commands}
	\begin{lstlisting}[caption={Aktivieren des Zeitplanungsprogrammes.}]
		elbrus@server:/var/elbrus$ sudo systemctl enable ssh-manager-schedule.timer
	\end{lstlisting}
	
	\lstset{style=commands}
	\begin{lstlisting}[caption={Starten des Zeitplanungsprogrammes.}]
		elbrus@server:/var/elbrus$ sudo systemctl start ssh-manager-schedule.timer
	\end{lstlisting}
	\newpage
	
	% --------- GEO SESSION FINDER ---------
	\section{Geo Session finder}
	% --------- deployment guide for the geo session finder ---------	
	\subsection{1 - Mit Setup Script}
	
	\lstset{style=commands}
	\begin{lstlisting}[caption={Ausführen des 'install.sh' Scripts.}]
		elbrus@server:~$ cd /var/elbrus
		elbrus@server:~$ sudo bash geo-session-finder/src/install.sh
		Do you want to proceed with setup of the 'geo session finder'? (y/n) y
		
		Where should the log be stored (dir) [/var/elbrus/shared/log]:
		Where should the 'geo-session-finder' be stored [/var/elbrus/geo-session-finder]:
		Where is the shared config stored [/var/elbrus/shared/.config]:
		
		Should the log be stored at '/var/elbrus/shared/log' ?
		Should the 'geo session finder' be stored at '/var/elbrus/geo-session-finder'?
		Is the shared config stored at '/var/elbrus/shared/.config' ? (y/n/exit) y
		
		<@\textcolor{codegreen}{Success!}@> .env file was created
		<@\textcolor{codegreen}{Success!}@> systemd service was automatically deployed,
		
		Installing dependencies ...
		
		...
		
		elbrus@server:~$	
	\end{lstlisting}

	\subsection{2 - Ohne Setup Script}

	\lstset{style=commands}
	\begin{lstlisting}[caption={Installieren der Abhängigkeiten.}]
		elbrus@server:~$ pip3 install -r geo_session_finder/requirements.txt
	\end{lstlisting}

	\subsubsection[file config]{Umgebung Konfigurieren}
	\lstset{style=files}
	\begin{lstlisting}[caption={Anhand von '.env.example' eigene '.env' Datei anlegen.}, language=bash]
		#global
		SHAREDCONFIG=/var/elbrus/shared/.config
		
		#paths
		LOGFILEDIR=/var/elbrus/shared/log
	\end{lstlisting}
	\newpage
	
	\subsubsection[systemd service]{Der Systemd Service}
	
	\lstset{style=files}
	\begin{lstlisting}[caption={geo-session-finder.service.example - Die Variable 'WorkingDirectory' sowie die Variable 'User' anpassen.},language=bash ,keywords={WorkingDirectory, User}, keywordstyle=\color{red}, firstnumber=5]
		...
		#job is starting immediatly after the start action has been called
		Type=simple
		#the user to execute the script
		User=elbrus
		#the working directory
		WorkingDirectory=/var/elbrus/geo_session_finder/src
		#which script should be executed
		ExecStart=/bin/bash elb-geo-session-finder.sh
		...
	\end{lstlisting}
	
	\lstset{style=commands}
	\begin{lstlisting}[caption={Kopieren des Serviceprogrammes.}]
		elbrus@server:/var/elbrus$ sudo cp geo_session_finder/src/geo-session-finder\
		.service.example /etc/systemd/system/geo-session-finder.service
	\end{lstlisting}
	
	\lstset{style=commands}
	\begin{lstlisting}[caption={Kopieren des Zeitplanungsprogrammes.}]
		elbrus@server:/var/elbrus$ sudo cp geo_session_finder/src/geo-session-finder-schedule\
		.timer.example /etc/systemd/system/geo-session-finder-schedule.timer
	\end{lstlisting}
	
	\lstset{style=commands}
	\begin{lstlisting}[caption={Neuladen des 'systemctl' Deamons.}]
		elbrus@server:/var/elbrus$ sudo systemctl daemon-reload
	\end{lstlisting}
	
	\lstset{style=commands}
	\begin{lstlisting}[caption={Aktivieren des Serviceprogrammes.}]
		elbrus@server:/var/elbrus$ sudo systemctl enable geo-session-finder.service
	\end{lstlisting}
	
	\lstset{style=commands}
	\begin{lstlisting}[caption={Aktivieren des Zeitplanungsprogrammes.}]
		elbrus@server:/var/elbrus$ sudo systemctl enable geo-session-finder-schedule.timer
	\end{lstlisting}
	
	\lstset{style=commands}
	\begin{lstlisting}[caption={Starten des Zeitplanungsprogrammes.}]
		elbrus@server:/var/elbrus$ sudo systemctl start geo-session-finder-schedule.timer
	\end{lstlisting}
	\newpage
	
	% --------- UPTIME MONITOR ---------
	\section{Uptime Monitor}
	% --------- deployment guide for the uptime monitor ---------	
	\subsection[file config]{Umgebung Konfigurieren}
	\subsubsection{1 - Mit Setup Script}
	
	\lstset{style=commands}
	\begin{lstlisting}[caption={Ausführen des 'install.sh' Scripts.}]
		elbrus@server:~$ cd /var/elbrus
		elbrus@server:~$ v
		Do you want to proceed with the setup of the 'uptime-monitor'? (y/n) y
		
		Where is the shared config stored [/var/elbrus/shared/.config]:
		Where should the 'uptime-monitor' be stored [/var/elbrus/uptime-monitor]:
		
		Is the shared config stored at '/var/elbrus/shared/.config'?
		Should the 'uptime-monitor be stored at '/var/elbrus/uptime-monitor'? (y/n/exit) y
		
		<@\textcolor{codegreen}{Success!}@> .env file was created
		<@\textcolor{codegreen}{Success!}@> systemd service was automatically deployed,
		elbrus@server:~$
	\end{lstlisting}
	
	\subsubsection{2 - Ohne Setup Script}
	
	\lstset{style=files}
	\begin{lstlisting}[caption={Anhand von '.env.example' eigene '.env' Datei anlegen.}, language=bash]
		#global
		SHAREDCONFIG=/var/elbrus/shared/.config
		
		#config
		# Initial pings to see if device is alive
		INITIALPING=1
		# Pings to get the availability statistic
		STATISTICPING=10
	\end{lstlisting}
	\newpage
	
	\subsection[systemd service]{Der Systemd Service}
	
	\lstset{style=files}
	\begin{lstlisting}[caption={uptime\_monitor.service.example - Die Variable 'WorkingDirectory' sowie die Variable 'User' anpassen.},language=bash ,keywords={WorkingDirectory, User}, keywordstyle=\color{red}, firstnumber=5]
		...
		#job is starting immediatly after the start action has been called
		Type=simple
		#the user to execute the script
		User=elbrus
		#the working directory
		WorkingDirectory=/var/elbrus/uptime_monitor
		#which script should be executed
		ExecStart=/bin/bash uptime_monitor.sh
		...
	\end{lstlisting}
	
	\lstset{style=commands}
	\begin{lstlisting}[caption={Kopieren des Serviceprogrammes.}]
		elbrus@server:/var/elbrus$ sudo cp uptime_monitor/uptime_monitor.service.example \
		/etc/systemd/system/uptime_monitor.service
	\end{lstlisting}
	
	\lstset{style=commands}
	\begin{lstlisting}[caption={Kopieren des Zeitplanungsprogrammes.}]
		elbrus@server:/var/elbrus$ sudo cp uptime_monitor/uptime_monitor-schedule.timer.example \
		/etc/systemd/system/uptime_monitor-schedule.timer
	\end{lstlisting}
	
	\lstset{style=commands}
	\begin{lstlisting}[caption={Neuladen des 'systemctl' Deamons.}]
		elbrus@server:/var/elbrus$ sudo systemctl daemon-reload
	\end{lstlisting}
	
	\lstset{style=commands}
	\begin{lstlisting}[caption={Aktivieren des Serviceprogrammes.}]
		elbrus@server:/var/elbrus$ sudo systemctl enable uptime_monitor.service
	\end{lstlisting}
	
	\lstset{style=commands}
	\begin{lstlisting}[caption={Aktivieren des Zeitplanungsprogrammes.}]
		elbrus@server:/var/elbrus$ sudo systemctl enable uptime_monitor-schedule.timer
	\end{lstlisting}
	
	\lstset{style=commands}
	\begin{lstlisting}[caption={Starten des Zeitplanungsprogrammes.}]
		elbrus@server:/var/elbrus$ sudo systemctl start uptime_monitor-schedule.timer
	\end{lstlisting}
	\newpage
	
	% --------- office365 ---------
	\section{office365}
	% --------- deployment guide for the office 365 get endpoint ---------
	\subsection[file config]{Umgebung Konfigurieren}
	\subsubsection{1 - Mit Setup Script}
	
	\lstset{style=commands}
	\begin{lstlisting}[caption={Ausführen des 'install.sh' Scripts.}]
		elbrus@server:~$ cd /var/elbrus
		elbrus@server:~$ sudo bash office365-analyzer/src/install.sh
		Do you want to proceed with setup of the 'office365-analyzer'? (y/n) y
		
		Where should the 'office365-analyzer' be stored (dir) [/var/elbrus/office365-analyzer]:
		Where should the log be stored (dir) [/var/elbrus/shared/log]:
		Where is the shared config stored [/var/elbrus/shared/.config]:
		
		Should the 'office365-analyzer' be stored at '/var/elbrus/office365-analyzer' ?
		Should the log be stored at '/var/elbrus/shared/log' ?
		Is the shared config stored at '/var/elbrus/shared/.config' ? (y/n/exit) y
		
		<@\textcolor{codegreen}{Success!}@> .env file was created
		<@\textcolor{codegreen}{Success!}@> systemd service was automatically deployed,
		
		info: installing dependencies.
		
		...
		
		info: installed dependencies.
		
		elbrus@server:~$	
	\end{lstlisting}
	
	\subsubsection{2 - Ohne Setup Script}
	
	\lstset{style=files}
	\begin{lstlisting}[caption={Anhand von '.env.example' eigene '.env' Datei anlegen.}, language=bash]
		#global
		SHAREDCONFIG=/var/elbrus/shared/.config
		
		#paths
		LOGFILEDIR=/var/elbrus/shared/log
		
		#ms url
		MS_URL=https://endpoints.office.com/endpoints/worldwide?clientrequestid=b10c5ed1-bad1-445f-b386-b919946339a7
	\end{lstlisting}
	\newpage

	\subsection[systemd service]{Der Systemd Service}
	
	\lstset{style=files}
	\begin{lstlisting}[caption={uptime\_monitor.service.example - Die Variable 'WorkingDirectory' sowie die Variable 'User' anpassen.},language=bash ,keywords={WorkingDirectory, User}, keywordstyle=\color{red}, firstnumber=4]
		...
		#job is starting immediatly after the start action has been called
		Type=simple
		#the user to execute the script
		User=elbrus
		#the working directory
		WorkingDirectory=/var/elbrus/office365-analyzer/src
		#which script should be executed
		ExecStart=/bin/bash elb-office365-get-endpoints.sh
		...
	\end{lstlisting}
	
	\lstset{style=commands}
	\begin{lstlisting}[caption={Kopieren des Serviceprogrammes.}]
		elbrus@server:/var/elbrus$ sudo cp office365-analyzer/src/office365-get-endpoints\
		.service.example /etc/systemd/system/office365-get-endpoints.service
	\end{lstlisting}
	
	\lstset{style=commands}
	\begin{lstlisting}[caption={Kopieren des Zeitplanungsprogrammes.}]
		elbrus@server:/var/elbrus$ sudo cp office365-analyzer/src/office365-get-endpoints\
		-schedule.timer.example /etc/systemd/system/office365-get-endpoints-schedule.timer
	\end{lstlisting}
	
	\lstset{style=commands}
	\begin{lstlisting}[caption={Neuladen des 'systemctl' Deamons.}]
		elbrus@server:/var/elbrus$ sudo systemctl daemon-reload
	\end{lstlisting}
	
	\lstset{style=commands}
	\begin{lstlisting}[caption={Aktivieren des Serviceprogrammes.}]
		elbrus@server:/var/elbrus$ sudo systemctl enable office365-get-endpoints.service
	\end{lstlisting}
	
	\lstset{style=commands}
	\begin{lstlisting}[caption={Aktivieren des Zeitplanungsprogrammes.}]
		elbrus@server:/var/elbrus$ sudo systemctl enable office365-get-endpoints-schedule.timer
	\end{lstlisting}
	
	\lstset{style=commands}
	\begin{lstlisting}[caption={Starten des Serviceprogrammes \& Zeitplanungsprogrammes.}]
		elbrus@server:/var/elbrus$ sudo systemctl start office365-get-endpoints.service
		elbrus@server:/var/elbrus$ sudo systemctl start office365-get-endpoints-schedule.timer
	\end{lstlisting}
	\newpage
	
	% --------- API ---------
	\section{API}
	% --------- deployment guide for the API ---------		
	\subsection{Voraussetzungen}
	
	\lstset{style=commands}
	\begin{lstlisting}[caption={Installieren von 'pm2'.}]
		elbrus@server:~/var/elbrus$ sudo npm install -g pm2
	\end{lstlisting}
	
	\lstset{style=commands}
	\begin{lstlisting}[caption={Nachinstallieren der Abhängigkeiten.}]
		elbrus@server:~/var/elbrus$ cd api
		elbrus@server:~/var/elbrus/api$ sudo npm install
	\end{lstlisting}
	\newpage
	
	\subsection[file config]{Umgebung Konfigurieren}
	
	\lstset{style=files}
	\begin{lstlisting}[caption={Anhand von '.env.example' eigene '.env' Datei anlegen.}, language=bash]
		# Application Name
		APP_NAME=Elbrus-API
		
		# Port number
		PORT=3000
		
		# BASE URL
		BASE=https://localhost:3000
		
		# URL of DB
		DB_USER=
		DB_HOST=
		DB_DATABASE=
		DB_PASSWORD=
		DB_PORT=
		
		# JWT
		JWT_SECRET=thisisasamplesecret
		JWT_ACCESS_EXPIRATION_MINUTES=30
		JWT_REFRESH_EXPIRATION_DAYS=30
		
		# SMTP configuration options for the email service
		SMTP_HOST=
		SMTP_PORT=
		SMTP_USERNAME=
		SMTP_PASSWORD=
		EMAIL_FROM=
		EMAIL_NAME=
	\end{lstlisting}
	
	\subsection{Inbetriebnahme}
	
	\lstset{style=commands}
	\begin{lstlisting}[caption={Starten der API.}]
		elbrus@server:~/var/elbrus/api$ pm2 start ecosystem.config.json
	\end{lstlisting}
	Die API läuft in folge automatisch im Hintergrund.
	\newpage
	
	\begin{enumerate}
		\item \textbf{APP\_NAME} wird rein als beschreibender Name genutzt und kann so belassen werden.
		\item \textbf{PORT} beschreibt den TCP Port auf dem die Applikation laufen soll.
		\item \textbf{BASE} ist der Wert der Basis URL auf welche zugegriffen wird. Hier muss der Port auch angegeben werden!
		\item \textbf{DB\_USER} ist der benutzername des DBMS Benutzers, über welchen der Zugriff auf die Datenbank läuft.
		\item \textbf{DB\_HOST} ist der hostname/ip-adresse des Servers welcher die Datenbank hostet.
		\item \textbf{DB\_DATABASE} beschreibt den Namen der Datenbank selber.
		\item \textbf{DB\_PASSWORD} ist das Passwort des DBMS Benutzers, über welchen der Zugriff auf die Datenbank läuft.
		\item \textbf{DB\_PORT} ist der TCP Port des Servers welcher die Datenbank hostet.
		\item \textbf{JWT\_SECRET} ist das Passwort mit dem alle JWT Tokens ausgestellt werden.
		\item \textbf{JWT\_ACCESS\_EXPIRATION\_MINUTES} gibt die Dauer der Gültigkeit eines Access-Tokens an (in Minuten)
		\item \textbf{JWT\_REFRESH\_EXPIRATION\_DAYS} gibt die Dauer der Gültigkeit eines Refresh-Tokens an (in Tagen)
		\item \textbf{SMTP\_HOST} ist der hostname/ip-adresse des EMail Servers
		\item \textbf{SMTP\_PORT} ist der TCP Port des EMail Servers für SMTP
		\item \textbf{SMTP\_USERNAME} ist der username des Benutzers zum einloggen in den EMail Account
		\item \textbf{SMTP\_PASSWORD} ist das passwort des Benutzers zum einloggen in den EMail Account
		\item \textbf{EMAIL\_FROM} gibt die Email adresse an, von welcher gesendet werden soll.
		\item \textbf{EMAIL\_NAME} gibt den Namen an, welcher dem Empfänger angezeigt werden soll.
	\end{enumerate}
	\newpage
	
	% --------- Webinterface ---------
	\section{Webinterface}
	% --------- deployment guide for the webinterface ---------		
	\subsection{1 - Mit Setup Script}
	
	\lstset{style=commands}
	\begin{lstlisting}[caption={Ausführen des 'install.sh' Scripts.}]
		elbrus@server:~$ cd /var/elbrus
		elbrus@server:/var/elbrus$ sudo bash webinterface/install.sh
			Do you want to proceed with the setup of the 'webinterface'? (y/n) y
			
			Where should the 'webinterface' be stored [/var/elbrus/webinterface]:
			
			Do you want to store the 'webinterface' at '/var/elbrus/webinterface'? (y/n/exit)
			 y
			
			info: installing dependencies:
			
			info: installing nginx
			info: installed nginx
			
			What is your baseurl (eg. 'http://1.2.3.4:3000/v1/' )? http://1.2.3.4/3000/v1/
			
			Is your baseurl 'http://1.2.3.4/3000/v1/'? (y/n/exit) y
			
			
			Success! .env file was created.
			
			
			info: building frontend
			
			> elbrus-webinterface@0.1.0 build /var/elbrus/webinterface
			> vue-cli-service build
			
			
			.:  Building for production...
			
			...
			
			info: built frontend
			info: overwriting nginx config. Your previous nginx config is stored as
			'nginx.conf.be'.
			info: starting nginx.
			Success: webinterface is now running at 'http://your-ip:80/'.
		elbrus@server:/var/elbrus$
	\end{lstlisting}
	\newpage
	
	\subsection{2 - Ohne Setup Script}
	
	\subsubsection[file config]{Umgebung Konfigurieren}
	\lstset{style=files}
	\begin{lstlisting}[caption={Anhand von '.env.example' eigene '.env' Datei anlegen.}, language=bash, keywords={CONFIGPATH, MAINPATH}, keywordstyle=\color{red}]
		VUE_APP_BASEURL=http://localhost:3000/v1/
	\end{lstlisting}
	
	\lstset{style=commands}
	\begin{lstlisting}[caption={Compilieren des Webinterface.}]
		elbrus@server:~$ cd /var/elbrus/webinterface
		elbrus@server:/var/elbrus/webinterface$ sudo npm run build
	\end{lstlisting}

	\lstset{style=commands}
	\begin{lstlisting}[caption={Kopieren des kompilierten Webinterfaces in der Ordner des Webservers.}]
		elbrus@server:/var/elbrus/webinterface$ sudo cp -r dist/ /usr/share/nginx/html/
	\end{lstlisting}

	\lstset{style=files}
	\begin{lstlisting}[caption={/etc/nginx/nginx.conf - Ändern der folgenden Zeilen.}, numbers=none, keepspaces=true]
		<@\textcolor{red}{	- server \{ }@>
		<@\textcolor{red}{	- 	listen       80 default\_server; }@>
		<@\textcolor{red}{	-	listen       [::]:80 default\_server; }@>
		<@\textcolor{red}{	-	server\_name  \_; }@>
		<@\textcolor{red}{	-	root         /usr/share/nginx/html; }@>
		<@\textcolor{red}{	-	\# Load configuration files for the default server block. }@>
		<@\textcolor{red}{	-	include /etc/nginx/default.d/*.conf; }@>
		<@\textcolor{red}{	-	location / \{ }@>
		<@\textcolor{red}{	- \} }@>
		
		<@\textcolor{codegreen}{ + server \{ }@>
		<@\textcolor{codegreen}{ +		listen       80 default\_server; }@>
		<@\textcolor{codegreen}{ +		listen       [::]:80 default\_server; }@>
		<@\textcolor{codegreen}{ +		server\_name  \_; }@>
		<@\textcolor{codegreen}{ +		root         /usr/share/nginx/html/dist; }@>
		<@\textcolor{codegreen}{ +		\# Load configuration files for the default server block. }@>
		<@\textcolor{codegreen}{ +		include /etc/nginx/default.d/*.conf; }@>
		<@\textcolor{codegreen}{ +		location / \{ }@>
		<@\textcolor{codegreen}{ +	\} }@>
		
		<@\textcolor{red}{	-	error\_page 404 \/404.html; }@>
		<@\textcolor{red}{	-			location = \/40x.html \{ }@>
		<@\textcolor{red}{	-	\} }@>
		
		<@\textcolor{codegreen}{ + error\_page 404 = @elbrus; }@>
		<@\textcolor{codegreen}{ + location @elbrus \{ }@>
		<@\textcolor{codegreen}{ +	root  /usr/share/nginx/html/dist;}@>
		<@\textcolor{codegreen}{ +	try\_files \$uri /index.html =502; }@>
		<@\textcolor{codegreen}{ + \} }@>
	\end{lstlisting}
	\newpage

	\subsection[systemd service]{Der Systemd Service}
	
	\lstset{style=commands}
	\begin{lstlisting}[caption={Aktivieren des Webservers.}]
		elbrus@server:/var/elbrus/webinterface$ sudo systemctl enable nginx
	\end{lstlisting}
	
	\lstset{style=commands}
	\begin{lstlisting}[caption={Starten des Webservers.}]
		elbrus@server:/var/elbrus/webinterface$ sudo systemctl systemctl start nginx 
	\end{lstlisting}
	\newpage
	\textcolor{red}{Achtung!} Vergessen sie nicht in ihren Firewall Einstellungen die Ports für die API sowie das Webinteface frei zu geben.
\end{document}